\section{Results}
\label{sec:results}
\noindent\textbf{This section is intentionally unpopulated.} Tables and figure slots below are the required structure for measured Kaggle outputs. Dashes indicate \emph{not yet measured}, not a zero score or a nonsignificant finding. The final text must be written after the dataset label audit and evaluation registry pass their checks. No quantitative outcomes are inferred from the software implementation.

\subsection{Real-data eligibility and frozen partitions}
\begin{table}[ht]
\centering\small
\caption{Final diagnosis-reviewed real-image inventory. Replace every dash with recorded reviewed counts and state the number of protective groups.}
\label{tab:realdata}
\begin{tabular}{@{}lrrrr@{}}
\toprule
Target class & Approved & Train & Val & Real test\\\midrule
Acne & \NA & \NA & \NA & \NA\\
Fungal infection & \NA & \NA & \NA & \NA\\
Vitiligo & \NA & \NA & \NA & \NA\\
Hyperpigmentation & \NA & \NA & \NA & \NA\\
Skin cancer/precancer scope & \NA & \NA & \NA & \NA\\
Normal & \NA & \NA & \NA & \NA\\\midrule
Total & \NA & \NA & \NA & \NA\\
\bottomrule
\end{tabular}
\end{table}
\noindent\Pending. Insert exclusion counts, exact/near-duplicate groups, cross-class conflicts, group-origin information and source-rights decisions. Explain if known patient IDs were unavailable.

\subsection{Generation yield and morphology review}
\begin{table}[ht]
\centering\small
\caption{Generator-prompt yield, accepted appearance quality and rejected duplicates. Each row must identify the actual generator revision and adaptation mode.}
\label{tab:synthetic}
\begin{tabularx}{\linewidth}{@{}Xrrrr@{}}
\toprule
Configuration & Candidates & Accepted & Accept (\%) & Near-copy rejects\\\midrule
E1 generic prompts & \NA & \NA & \NA & \NA\\
E2 morphology prompts & \NA & \NA & \NA & \NA\\
Optional LoRA source (separately) & \NA & \NA & \NA & \NA\\
\bottomrule
\end{tabularx}
\end{table}
\noindent\Pending. Insert a permitted representative-image grid only after originality, privacy, attribution and morphology review. The figure must distinguish real and synthetic images unambiguously.

\subsection{Primary real-test classification}
\begin{table}[ht]
\centering\small
\caption{Primary real-test performance by experiment arm. Make one version per backbone or add a backbone column when measured. Report mean across seeds and its SD separately.}
\label{tab:primary}
\begin{tabular}{@{}lrrr@{}}
\toprule
Arm & Macro-F1 & Balanced accuracy & Accuracy\\\midrule
E0 real only & \NA & \NA & \NA\\
E1 generic synthetic & \NA & \NA & \NA\\
E2 morphology synthetic & \NA & \NA & \NA\\
E3 quality selected & \NA & \NA & \NA\\
E4 diversity selected (optional) & \NA & \NA & \NA\\
E5 synthetic only (optional) & \NA & \NA & \NA\\
\bottomrule
\end{tabular}
\end{table}
\noindent\Pending. State the tested backbone, exact frozen real-test sample size, number of distinct protective groups, matched synthetic budget, training seeds, checkpoint version and failure/early-stopping events.

\subsection{Matched real-test contrasts}
\begin{table}[ht]
\centering\small
\caption{Paired differences in real-test macro-F1, evaluated on the identical real test for each backbone and seed. A CI is not a p-value.}
\label{tab:paired}
\begin{tabularx}{\linewidth}{@{}lXr@{}}
\toprule
Contrast & Interpretation & Difference [95\% group-bootstrap CI]\\\midrule
E2 versus E0 & Augmentation versus real only & \NA\\
E2 versus E1 & Morphological specificity versus generic prompts & \NA\\
E3 versus E2 & Quality selection at matched sample count & \NA\\
E4 versus E2 & Diversity selection at matched sample count & \NA\\
\bottomrule
\end{tabularx}
\end{table}
\noindent\Pending. Insert paired bootstrap intervals, observed signs across training seeds and a statement on the reliability of the underlying grouping unit. Do not add a fabricated significance column.

\subsection{Class errors, source behavior and calibration}
\noindent\Pending. Insert real-test six-class confusion matrices, per-class precision/recall/F1, source-specific support and calibrated probability diagnostics. Figures should come from the saved Evaluation 04--07 artifacts, with sources shown only at justified support thresholds. If synthetic prompt-test results exist, place them in a separate supplementary panel rather than pooling them into the primary real-test score.

\subsection{Result-dependent synthesis}
\noindent\Pending. Replace this single paragraph with a factual summary answering RQ1--RQ3, naming the backbones, effect magnitudes, uncertainty intervals, disagreement across seeds, and negative results. A clinically favorable interpretation must be justified by the data and cannot rest on visual quality alone.
